\documentclass[inzynier,druk,ramka]{dyplom}
\usepackage[utf8]{inputenc}
\usepackage{hyperref}
%%
\usepackage[toc]{appendix}
\renewcommand{\appendixtocname}{Dodatki}
\renewcommand{\appendixpagename}{Dodatki}
%%
\usepackage{listings}

%\noappendicestocpagenum
\usepackage{lipsum}
\usepackage{url}
%\usepackage{showframe}


\author{Jan A. Backi}
\title{Lorem ipsum dolor sit amet, consectetuer adipiscing elit}
\titlen{Lorem ipsum dolor sit amet, consectetuer adipiscing elit}
\promotor{dr hab. inż. Jerzy Babacki, prof. nadzw. PWr., I-77}
%\konsultant{dr hab. inż. Kazimerz Kabacki}
\wydzial{Wydział Mechaniczny}
\kierunek{Prawo}
\specjalnosc{Lewo}

\begin{document}

\maketitle

\tableofcontents

\listoffigures

\listoftables

%Ala ma kota

% !TeX program = latexmk
% !TeX spellcheck = pl_PL
% !TeX root = example.tex

\chapter{Krótka instrukcja obsługi}

\section{Wstęp}

Klasa przygotowana jest zgodnie z zaleceniami dostępnymi ze strony \url{} i może być wykorzystana do składu pracy \textbf{inżynierskiej} lub \textbf{magisterskiej} na wydziale mechanicznym.

Klasa zgodnie z \href{http://www.wmech.pwr.wroc.pl/88428.dhtml}{wymaganiami Wydziału Mechanicznego} składa stronę tytułową i~stosuje się do zaleceń (czcionka, zasady numeracji, odstępy,\ldots).

\section{Użycie}

\begin{enumerate}
\item
Praca magisterska i inżynierska.

Wychodzi, że tak na prawdę powinny być dwie wersje pracy: do archiwum (marginesy 2,5 cm) i „dla promotora”\footnote{Ciekawe po co mu…?}. Wersja dla promotora powinna mieć trochę większy margines od strony oprawy (35mm), najprawdopodobniej będzie drukowana \textbf{jednostronnie} i, żeby była łatwiejsza do czytania będzie złożona z interlinią \textbf{1,5}.

Jak tak to tak:  pojawiły się dwie dodatkowe opcje klasy:
\begin{enumerate}
\item
\texttt{archiwum}: \verb|\usepackage[magister,archiwum]{dyplom}| — wersja do archiwum
\item
\texttt{druk}: \verb|\usepackage[inzynier,druk]{dyplom}| — wersja do „druku” (i oprawy).
\end{enumerate}
\textbf{W przypadku braku opcji — wybierana jest wersja do archiwum!}

Tak na prawdę, to w przypdku braku opcji powinna być wybierana wersja druk. Wybrałem jednak opisane wyzej zachowanie, aby zachowanie zmodyfikowanej klasy było zgodne z dotychczasowym. Zalecam przeprowadzanie redakcji tekstu w trybie druku i pozostawienie dokumentu „jak wyjdzie” w trybie archiwum. Chyba, że ilustracje będą zachowywać się bardzo dziwnie…

Ponieważ „doszły do mnie” jakieś dziwne informacje, że ze stroną tytułową jest coś nie tak, dokonałem kolejnych porównań. Różnica jest jedna: obecność ramki wokół tytułów pracy. W związku z tym, ramka została zlikwidowana. Można ją odzyskać dodając dodatkowy parametr: \verb|\usepackage[inzynier,druk,ramka]{dyplom}| i~się pojawi…
\item
Praca magisterska:
\begin{verbatim}
\documentclass[magister]{dyplom}
\end{verbatim}
Dodatkowo zdefiniować należy sposób kodowania polskich liter. W przypadku systemu Windows będzie to najprawdopodobniej:
\begin{verbatim}
\usepackage[cp1250]{inputenc}
\end{verbatim}
a w przypadku systemów linuksowych:
\begin{verbatim}
\usepackage[utf8]{inputenc}
\end{verbatim}

Dodatkowo zdefiniować należy „metadane”:
\begin{itemize}
\item
Nazwisko autora:
\begin{verbatim}
\author{Imię Nazwisko}
\end{verbatim}
\item
Tytuł pracy (w języku polskim):
\begin{verbatim}
\title{Tytuł Pracy}
\end{verbatim}
\item
Tytuł pracy po angielsku
\begin{verbatim}
\titlen{Work Title}
\end{verbatim}
\item
Nazwisko promotora
\begin{verbatim}
\promotor{prof. dr hab. inż. Imię Nazwisko, prof. PWr.}
\end{verbatim}
\item
Kierunek
\begin{verbatim}
\kierunek{Prawo}
\end{verbatim}
\item
Specjalność:
\begin{verbatim}
\specjalnosc{Lewo}
\end{verbatim}
\item
W razię potrzeby wpisać można inną nazwę wydziału. Gdy nie zostanie wpisana — będzie tam Wydział Mechaniczny.
\begin{verbatim}
\wydzial{Wydział Elektryczny}
\end{verbatim}
\item
Praca może mieć konsultanta/konsultantów. Dodałem więc pole konsultant:
\begin{verbatim}
\konsultant{dr inż. Kazimierz Kabacki}
\end{verbatim}
Nazwisko konsultanta pojawi się miedzy nazwiskiem promotora a oceną. Pozostaje kwestia czy powinien to być „konsultant” czy raczej „konsulktanci”?
\end{itemize}
Powyższe metadane umieszczamy przed \verb|\begin{document}|:
\begin{verbatim}
\documentclass[magister]{dyplom}
\usepackage[utf8]{inputenc}

\author{Jan A. Backi}
\title{Lorem ipsum dolor sit amet, consectetuer adipiscing elit}
\titlen{Lorem ipsum dolor sit amet, consectetuer adipiscing elit}
\promotor{dr hab. inż. Jerzy Babacki, prof. nadzw. PWr., I-77}
\wydzial{Wydział Mechaniczny}
\kierunek{Prawo}
\specjalnosc{Lewo}

\begin{document}
\end{verbatim}
\item
Praca inżynierska:
\begin{verbatim}
\documentclass[inzynier]{dyplom}
\end{verbatim}
Dodatkowo zdefiniować należy sposób kodowania polskich liter. W przypadku systemu Windows będzie to najprawdopodobniej:
\begin{verbatim}
\usepackage[cp1250]{inputenc}
\end{verbatim}
a w przypadku systemów linuksowych:
\begin{verbatim}
\usepackage[utf8]{inputenc}
\end{verbatim}

Dodatkowo zdefiniować należy „metadane”:
\begin{itemize}
\item
Nazwisko autora:
\begin{verbatim}
\author{Imię Nazwisko}
\end{verbatim}
\item
Tytuł pracy (w języku polskim):
\begin{verbatim}
\title{Tytuł Pracy}
\end{verbatim}
\item
Tytuł pracy po angielsku
\begin{verbatim}
\titlen{Work Title}
\end{verbatim}
\item
Nazwisko promotora
\begin{verbatim}
\promotor{prof. dr hab. inż. Imię Nazwisko, prof. PWr.}
\end{verbatim}
\item
Kierunek
\begin{verbatim}
\kierunek{Prawo}
\end{verbatim}
%\item
%Specjalność:
%\begin{verbatim}
%\specjalnosc{Lewo}
%\end{verbatim}
\item
W razię potrzeby wpisać można inną nazwę wydziału. Gdy nie zostanie wpisana — będzie tam Wydział Mechaniczny.
\begin{verbatim}
\wydzial{Wydział Elektryczny}
\end{verbatim}
\item
Praca może mieć konsultanta/konsultantów. Dodałem więc pole konsultant:
\begin{verbatim}
\konsultant{dr inż. Kazimierz Kabacki}
\end{verbatim}
Nazwisko konsultanta pojawi się miedzy nazwiskiem promotora a oceną. Pozostaje kwestia czy powinien to być „konsultant” czy raczej „konsultanci”?

\end{itemize}
Powyższe metadane umieszczamy przed \verb|\begin{document}|:
\begin{verbatim}
\documentclass[magister]{dyplom}
\usepackage[utf8]{inputenc}

\author{Jan A. Backi}
\title{Lorem ipsum dolor sit amet, consectetuer adipiscing elit}
\titlen{Lorem ipsum dolor sit amet, consectetuer adipiscing elit}
\promotor{dr hab. inż. Jerzy Babacki, prof. nadzw. PWr., I-77}
\wydzial{Wydział Mechaniczny}
\kierunek{Prawo}
\specjalnosc{Lewo}

\begin{document}
\end{verbatim}

\end{enumerate}

\section{Dodatkowe zasoby}

Warto wspomnieć  o innych inicjatywach przyswojenia LaTeX{}a piszącym prace dyplomowe. Najważniejsza z nich to książka Tomasza Przechlewskiego \cite{tp-11-latex} oraz przygotowana przez niego klasa znajdująca się w \url{https://github.com/hrpunio/wzmgr}. Przykłady z książki znaleźć można w \url{https://github.com/hrpunio/pmdzpl}.


\section{Gdzie znaleźć?}

Pakiet można znaleźć pod adresem: \url{http://www.immt.pwr.wroc.pl/~myszka/dydaktyka/}. Wersja zarchiwizowana: \href{http://www.immt.pwr.wroc.pl/~myszka/TeX/dyplom/dyplom.zip}{dyplom.zip}

\section{Uwagi}

Wszelkie uwagi i postulaty należy kierować na adres Wojciech.Myszka@pwr.wroc.pl

W miarę potrzeby mogę szablon dostosować do wymagań innych wydziałow Politechniki Wrocławskiej.



\chapter{Ala ma kota}

ĄĆĘŁŃÓŚŹŻ

ąćęłńóśźż

\lipsum[1]

\section{Ola ma psa}

\lipsum[2-3]

\section{Zosia ma kosia}

\lipsum[7]

\subsection{Ula ma małpę}

\lipsum[4-10]

\begin{figure}
\includegraphics[width=.4\textwidth]{kotek}
\caption{Ala ma kota}
\end{figure}

\lipsum[11-15]

\begin{table}
\caption{Co kto ma \cite{harel_rzecz_2008} (patrz też dodatek~\ref{Dod1})}
\begin{tabular}{|l|l|l|}
\hline
Ala & ma & kota \\
\hline
Ola & ma & psa \\
\hline
Ula & ma & małpę\\
\hline
\end{tabular}
\end{table}

\lipsum[16-20]

\begin{equation}
\sum_{i=1}^{\infty}a_i
\end{equation}

\chapter{Listingi}

\begin{lstlisting}
int main()
{
   int a=2*3;
   printf("**Ala ma kota\n**");
   while(!I2C_CheckEvent(I2C1, I2C_EVENT_MASTER_MODE_SELECT)); /* EV5 */
   return 0;
}
\end{lstlisting}


\appendixpage
\appendix
%\addappheadtotoc

\chapter{To powinien być dodatek}\label{Dod1}

\lipsum[9-11]

\nocite{*}
\bibliography{biblioex}
\bibliographystyle{dyplom}

\end{document}